\documentclass[12pt]{article}
\usepackage[utf8]{inputenc}
\usepackage{amsmath, amssymb, amsthm}
\usepackage{geometry}
\geometry{a4paper, margin=1in}

\newtheorem{theorem}{Theorem}
\newtheorem{lemma}[theorem]{Lemma}
\newtheorem{corollary}[theorem]{Corollary}
\newtheorem{proposition}[theorem]{Proposition}
\theoremstyle{definition}
\newtheorem{definition}[theorem]{Definition}

\title{Problem 53: Combinatoric Selections}
\author{Project Euler Solution}
\date{}

\begin{document}
\maketitle

\section{Problem Statement}
How many values of $\binom{n}{r}$ for $1 \leq n \leq 100$ and $0 \leq r \leq n$ are greater than one million ($10^6$)?

\section{Formal Development}

\begin{definition}[Binomial Coefficient]
For non-negative integers $n \geq r$, $\binom{n}{r} = \frac{n!}{r!(n-r)!}$.
\end{definition}

\begin{theorem}[Unimodality]
\label{thm:unimodal}
For fixed $n \geq 1$, the sequence $\binom{n}{0}, \binom{n}{1}, \ldots, \binom{n}{n}$ is strictly unimodal with maximum at $r = \lfloor n/2 \rfloor$.
\end{theorem}

\begin{proof}
The ratio $\binom{n}{r}/\binom{n}{r-1} = (n - r + 1)/r$ exceeds~1 if and only if $r < (n+1)/2$. Hence the sequence strictly increases for $r < (n+1)/2$ and strictly decreases for $r > (n+1)/2$.
\end{proof}

\begin{theorem}[Symmetry]
For all $0 \leq r \leq n$: $\binom{n}{r} = \binom{n}{n-r}$.
\end{theorem}

\begin{proof}
Immediate from the definition: $\binom{n}{n-r} = n!/((n-r)!\,r!) = \binom{n}{r}$.
\end{proof}

\begin{corollary}[Contiguous Exceedance]
\label{cor:contiguous}
For threshold $T > 0$ and fixed $n$, the set $E_T(n) = \{r : \binom{n}{r} > T\}$ is either empty or a contiguous interval $\{r_{\min}, \ldots, n - r_{\min}\}$ with $|E_T(n)| = n - 2r_{\min} + 1$.
\end{corollary}

\begin{proof}
Follows from unimodality (Theorem~\ref{thm:unimodal}) and symmetry: the exceedance set is a connected, symmetric sublevel set of a unimodal function.
\end{proof}

\begin{lemma}[Pascal's Recurrence]
For $1 \leq r \leq n - 1$: $\binom{n}{r} = \binom{n-1}{r-1} + \binom{n-1}{r}$.
\end{lemma}

\begin{proof}
Partition the $r$-element subsets of $\{1, \ldots, n\}$ by membership of element~$n$: those containing~$n$ number $\binom{n-1}{r-1}$; those not containing~$n$ number $\binom{n-1}{r}$.
\end{proof}

\begin{lemma}[Overflow Avoidance]
\label{lem:cap}
Define $\hat{C}(n, r) = \min(\binom{n}{r}, T+1)$. This satisfies the capped recurrence $\hat{C}(n, r) = \min(\hat{C}(n-1, r-1) + \hat{C}(n-1, r),\; T+1)$, and $\binom{n}{r} > T \iff \hat{C}(n,r) = T + 1$.
\end{lemma}

\begin{proof}
By induction. If both capped summands equal their true values (both $\leq T$), the sum is exact; if any summand is capped at $T+1$, the true sum exceeds $T$, and capping to $T+1$ preserves the exceedance property.
\end{proof}

\begin{theorem}[Counting Formula]
The total count is $\sum_{n=1}^{100} |E_{10^6}(n)|$.
\end{theorem}

\section{Algorithm}

\begin{enumerate}
    \item Initialize $\mathrm{prev} \gets [1]$ (row $n = 0$), $\mathrm{count} \gets 0$.
    \item For $n = 1$ to $100$:
    \begin{enumerate}
        \item Compute row $n$ via Pascal's recurrence with capping (Lemma~\ref{lem:cap}).
        \item Add $|\{r : \hat{C}(n, r) > T\}|$ to count.
    \end{enumerate}
    \item Return count.
\end{enumerate}

\section{Complexity}

\begin{proposition}[Time]
$O(N^2)$ where $N = 100$. Specifically, $\sum_{n=1}^{100}(n+1) = 5150$ operations.
\end{proposition}

\begin{proposition}[Space]
$O(N)$ for one row of Pascal's triangle (at most 101 entries).
\end{proposition}

\section{Result}
\[
\boxed{4075}
\]

\section{Answer}

\[
\boxed{4075}
\]

\end{document}
